\section{A rational trivial-sector obstruction with both polar conditions}
\label{sec:rational-physical-projection-obstruction}

\paragraph{Proof-review status.}
The complete author and independent whole-proof readings have
passed. Their exact proof and source-interface scopes are recorded
in the separate check events.

We now use exactly the rational trivial sector, $S=\{\infty\}$
and $S'=S$, with no added unit-shell transition places,
and the actual theta map and PW source of (MA-GP.1)--(MA-GP.2):
\[
 \begin{gathered}
 \mathcal U_0=\{f\in\mathcal S(\mathbb R): f(-t)=f(t),\quad
                         f(0)=0,\ \int_{\mathbb R}f(t)dt=0\},\\
 (\Theta f)(x)=2e^{x/2}\sum_{n\ge1}f(ne^x),\qquad
 \mathcal U_{0,Y}=\{f\in\mathcal U_0:
                  \operatorname{supp}\widehat f\subseteq[-Y,Y]\}.
 \end{gathered}
 \tag{MA-PI.70}
\]
The Fourier kernel is $e^{-2\pi it\xi}$. The range
$\Theta\mathcal U_0$ is in the all-exponential Schwartz
space $\mathcal E$, and the map into that space is continuous.
This is the actual source identification of MA-TR-101 and
the Poisson identity (MA-GP.8), not a prescribed divisor model.

Two specified Hilbert metrics will be treated independently:
\[
 \begin{aligned}
 H^{\rm sym}_{a,s}
  &=L^2\bigl(\mathbb R,e^{2a|x|}(1+x^2)^{s/2}dx\bigr),
             &&0<a<1/2,\ s\ge0,\\
 H^{\rm one}_{b,s}
  &=L^2\bigl(\mathbb R,e^{2bx}(1+x^2)^{s/2}dx\bigr),
             &&|b|<1/2,\ s\ge0.
 \end{aligned}
 \tag{MA-PI.71}
\]
In either one of these fixed metrics, put
\[
 J_Y=\overline{\Theta\mathcal U_{0,Y}},\qquad Q_Y=P_{J_Y}.
 \tag{MA-PI.72}
\]
Thus $Q_Y$ is its actual orthogonal projection. In the symmetric
case it is exactly the corrected global projector of MA-GP.
The one-sided assertion is a separately proved statement in
its displayed metric; equivalent norms are not used to transfer
orthogonal projections.

Let $\rho_Y\in C^\infty([0,\infty))$ satisfy
\[
 0\le\rho_Y\le1,\quad \rho_Y(r)=1\ (r\le Y/2),\quad
 \rho_Y(r)=0\ (r\ge2Y),\qquad
 \mathfrak b_Y=\mathcal F^{-1}M_{\rho_Y(|\xi|)}\mathcal F.
 \tag{MA-PI.73}
\]
This is the original physical Fourier cutoff for this $S$.

\paragraph{MA-PI-003 (the principal rational canonical intertwining fails).}
For either fixed metric in (MA-PI.71), there is $Y_0$ such that,
for every $Y\ge Y_0$ and every profile (MA-PI.73), there exists
an even Schwartz function $f_Y$, identically zero near the
additive origin, satisfying
\[
 f_Y\in\mathcal U_0,\qquad \mathfrak b_Y f_Y\in\mathcal U_0,
 \qquad Q_Y\Theta f_Y\ne\Theta\mathfrak b_Y f_Y.
 \tag{MA-PI.74}
\]
Both actual theta images therefore belong to $\mathcal E$.
The failure is not caused by an undefined polar output.
The input lies in the original flat source, with its two
global polar conditions. No information about the locations
or multiplicities of zeta zeros is used.

\begin{proof}
\emph{1. Enforcing the two source moments before the Gram calculation.}
Fix $b\in(-1/2,1/2)$ for this calculation and set
$\sigma=1/2+b$, $d=1-\sigma$.
Given $B\in C_c^\infty(\mathbb R)$, put
\[
 A=(\partial_r+\sigma)B,\qquad
 u(\xi)=\xi^{\sigma-1}A(\log\xi)\quad(\xi>0),
 \qquad u(-\xi)=u(\xi),\qquad f=\mathcal F^{-1}u.
 \tag{MA-PI.75}
\]
The frequency vector is smooth with compact support away from
zero. Integration by parts gives
\[
 \int_0^\infty u(\xi)d\xi
       =\int_{\mathbb R}e^{\sigma r}
                  (B'(r)+\sigma B(r))dr=0.
 \tag{MA-PI.76}
\]
Hence $f(0)=2\int_0^\infty u=0$ and
$\int f=\widehat f(0)=u(0)=0$. Thus $f\in\mathcal U_0$.
The differential operation does not enlarge the support of
$B$ in logarithmic frequency.

Write $M_f(z)=\int_0^\infty f(t)t^{z-1}dt$ and
$\widehat B(t)=\int B(r)e^{-itr}dr$. Actual unfolding, first
in $\Re z>1$ and then by continuation, gives
\[
 \mathcal M(\Theta f)(z-1/2)=2\zeta(z)M_f(z).
 \tag{MA-PI.77}
\]
The even inverse Fourier integral, with an Abel factor
$e^{-\varepsilon t}$ before interchanging integrals, gives
\[
 M_f(z)=\frac{\Gamma_{\mathbb R}(z)}
                    {\Gamma_{\mathbb R}(1-z)}M_u(1-z),
                         \qquad0<\Re z<1.
 \tag{MA-PI.78}
\]
Indeed $f(t)=2\int_0^\infty u(r)\cos(2\pi tr)dr$.
The Abel integral tends to
$2(2\pi)^{-z}\Gamma(z)\cos(\pi z/2)M_u(1-z)$,
uniformly over the compact positive support of $u$;
reflection and duplication give precisely (MA-PI.78).
On the line $z=\sigma+it$, one has
$M_u(1-z)=(\sigma+it)\widehat B(t)$.

The actual Poisson identity $\Theta f(x)=\Theta\widehat f(-x)$
in (MA-GP.8), together with (MA-PI.77)--(MA-PI.78), fixes
the functional-equation normalization
\[
 \frac{\Gamma_{\mathbb R}(z)}{\Gamma_{\mathbb R}(1-z)}
                  \zeta(z)=\zeta(1-z).
 \tag{MA-PI.79}
\]
For example use a nonzero $f$ constructed in (MA-PI.75)
in the common strip, cancel its nonzero Mellin transform on an open set,
and use meromorphic continuation. This is only the usual
rational functional equation in the precise Fourier convention
needed here; it contains no zero-location assertion.

For $B,C$ and their corresponding $f,h$, Mellin Plancherel now gives
\[
 \begin{aligned}
 \int_{\mathbb R}(\Theta f)(x)
       \overline{(\Theta h)(x)}e^{2bx}dx
    &=\frac1{2\pi}\int_{\mathbb R}
             \mathscr V_b(t)\widehat B(t)
                               \overline{\widehat C(t)}dt,\\
 \mathscr V_b(t)&=4(t^2+\sigma^2)
                       \zeta(d+it)\zeta(d-it).
 \end{aligned}
 \tag{MA-PI.80}
\]
In particular this is a Gram form of genuine sources with
both polar conditions, not a formal positive symbol.

\emph{2. The constrained Gram distribution has unbounded support.}
The real-axis function $\mathscr V_b$ is even, nonnegative,
and polynomially bounded. To check the last property directly,
the usual first partial-summation identity
\[
 \zeta(z)=\frac{z}{z-1}
               -z\int_1^\infty \{x\}x^{-z-1}dx,
                   \qquad\Re z>0,\ z\ne1,
 \tag{MA-PI.81}
\]
gives $\zeta(d+it)=O_d(1+|t|)$ on the fixed line
$0<d<1$. The integral follows initially from summing
$\lfloor x\rfloor$ in $\Re z>1$ and continues there
absolutely for $\Re z>0$. Thus $\mathscr V_b$ defines
a tempered distribution.

The same expression in a complex variable $t$ is entire.
Its only possible poles are the simple zeta poles at
$t=\pm i\sigma$; the factor $t^2+\sigma^2$ cancels
both. Its growth on the imaginary axis is not of
exponential type. From (MA-PI.79) and the gamma recurrence,
\[
 \zeta(-2n-1)=(-1)^{n+1}
       \frac{2(2n+1)!}{(2\pi)^{2n+2}}\zeta(2n+2),
                         \qquad n\ge0.
 \tag{MA-PI.82}
\]
Put $T_n=2n+1+d$. Since both zeta values at the positive
arguments in the following display exceed one, for $n\ge1$
\[
 \begin{aligned}
 |\mathscr V_b(iT_n)|
  &=4(T_n^2-\sigma^2)
          |\zeta(-2n-1)|\zeta(2n+1+2d)\\
  &\ge \frac{8(T_n^2-\sigma^2)(2n+1)!}{(2\pi)^{2n+2}}.
 \end{aligned}
 \tag{MA-PI.83}
\]
This exceeds every bound $C(1+T_n)^N e^{RT_n}$ eventually.
If the inverse Fourier distribution $\kappa_b$ of
$\mathscr V_b$ had compact support, its finite order on a
fixed compact neighborhood, applied to a cutoff times
$e^{-izx}$, would give exactly such an entire exponential
bound, as proved in (MA-PI.18). Hence $\kappa_b$ is not
compactly supported. It is even.

Consequently, for every two cut positions $A_0<B_0$ there
are compact smooth $B,C$, supported respectively in
$(B_0,\infty)$ and $(-\infty,A_0)$, with the Gram pairing
(MA-PI.80) nonzero. Here is the exact support argument.
If all such pairings vanished, testing $B$ by a smooth
approximation to a delta at $r_0>B_0$ would show that
$\kappa_b$ vanishes on $(-\infty,A_0-B_0)$. Each compact
test in that interval permits an $r_0$ with its translate
below $A_0$. Evenness would give vanishing beyond the
opposite end too, contradicting unbounded support.
A nonzero limiting distribution pairing already has a
nonzero ordinary smooth approximate-delta member. Thus
the resulting $B,C$, and hence $u,v,f,h$, are actual
smooth functions throughout.

Choose $A_0=\log(1/2)$ and $B_0=\log2$. Equations
(MA-PI.75)--(MA-PI.76) supply $f,h\in\mathcal U_0$ with
\[
 \operatorname{supp}\widehat f\subset\{|\xi|>2\},\qquad
 \operatorname{supp}\widehat h\subset\{0<|\xi|<1/2\},
 \quad D_b:=\langle\Theta f,\Theta h\rangle_{L^2(e^{2bx}dx)}
                              \ne0.
 \tag{MA-PI.84}
\]
These supports are also compact. No arbitrary preassigned
pair of bounded bands is asserted to work.

\emph{3. Physical dilation in each specified metric.}
Put $Y=e^L$, $f_L(t)=e^{L/2}f(e^Lt)$ and
$h_L(t)=e^{L/2}h(e^Lt)$. The two moments stay zero,
and actual summation gives
\[
 \Theta f_L(x)=\Theta f(x+L),\qquad
 \Theta h_L(x)=\Theta h(x+L),\qquad
 \mathfrak b_Y f_L=0,\qquad h_L\in\mathcal U_{0,Y}.
 \tag{MA-PI.85}
\]
The last two assertions use only the plateaus in
(MA-PI.73), so they hold uniformly over their profiles.
For the symmetric metric choose $b=-a$ in step 1.
Then the same dominated-convergence calculation as
(MA-PI.25) gives
\[
 e^{-2aL}L^{-s}
       \langle\Theta f_L,\Theta h_L\rangle_{H^{\rm sym}_{a,s}}
                        \longrightarrow D_{-a}\ne0.
 \tag{MA-PI.86}
\]
For the one-sided metric use its own parameter $b$ in
step 1. Direct substitution $x\mapsto x-L$ gives
\[
 e^{2bL}L^{-s}
       \langle\Theta f_L,\Theta h_L\rangle_{H^{\rm one}_{b,s}}
 =\int\Theta f(x)\overline{\Theta h(x)}e^{2bx}
       \frac{(1+(x-L)^2)^{s/2}}{L^s}dx
                         \longrightarrow D_b\ne0.
 \tag{MA-PI.87}
\]
For $L\ge1$ the polynomial ratio is at most
$C_s(2+|x|)^s$ and tends to one. Both theta images are
in $\mathcal E$, which proves domination also for $s=0$.
There is therefore one $Y_0$ for the fixed metric such
that the pairing is nonzero for all $Y\ge Y_0$,
independently of the profile. Since $\Theta h_L\in J_Y$,
\[
 \langle Q_Y\Theta f_L-\Theta\mathfrak b_Y f_L,
                      \Theta h_L\rangle_H
                 =\langle\Theta f_L,\Theta h_L\rangle_H\ne0.
 \tag{MA-PI.88}
\]
The bandlimited input $f_L$ is not flat. The final step
produces a genuine flat input and retains both polar
conditions also after the cutoff.

\emph{4. Two fixed flat correctors for a fixed physical profile.}
Fix one $Y\ge Y_0$ and its profile. The inverse Fourier
transform $k_Y=\mathcal F^{-1}(\rho_Y(|\xi|))$ is a real
even Schwartz function, and is not constant. Indeed a
constant Schwartz function is zero, whereas the Fourier
multiplier equals one near zero. On even compact smooth
functions supported away from the origin the two
functionals
\[
 I(v)=\int_{\mathbb R}v(t)dt,\qquad
 \ell_Y(v)=(\mathfrak b_Yv)(0)
                         =\int_{\mathbb R}k_Y(t)v(t)dt
 \tag{MA-PI.89}
\]
are linearly independent. Otherwise testing arbitrarily
small even bumps away from zero would force $k_Y$ to be
constant on $\mathbb R\setminus\{0\}$, and hence everywhere
by continuity. Thus there exist fixed real even functions
$\psi_1,\psi_2\in C_c^\infty(\mathbb R\setminus\{0\})$
such that
\[
 T_Y=\begin{pmatrix}
 I(\psi_1)&I(\psi_2)\\
 \ell_Y(\psi_1)&\ell_Y(\psi_2)
 \end{pmatrix}
                         \quad\hbox{is invertible}.
 \tag{MA-PI.90}
\]
This is finite-dimensional linear algebra after the
functional independence; no uniform choice over profiles
is claimed or needed.

Choose a fixed even $\omega\in C_c^\infty(\mathbb R)$
equal to one near zero. For $0<\varepsilon\le1$, put
\[
 \begin{gathered}
 r_\varepsilon(t)=\omega(t/\varepsilon)f_L(t),\qquad
 \binom{c_1(\varepsilon)}{c_2(\varepsilon)}
       =T_Y^{-1}\binom{I(r_\varepsilon)}{\ell_Y(r_\varepsilon)},\\
 f_{L,\varepsilon}=f_L-r_\varepsilon
                    +c_1(\varepsilon)\psi_1
                    +c_2(\varepsilon)\psi_2,\qquad
 e_\varepsilon=f_L-f_{L,\varepsilon}.
 \end{gathered}
 \tag{MA-PI.91}
\]
The function $f_{L,\varepsilon}$ is even, Schwartz, and
zero near zero. Equations (MA-PI.85) and (MA-PI.90)--(MA-PI.91)
give $I(f_{L,\varepsilon})=0$ and
$\ell_Y(f_{L,\varepsilon})=0$. Since the multiplier
is one at zero, also
$I(\mathfrak b_Y f_{L,\varepsilon})=I(f_{L,\varepsilon})=0$.
Consequently
\[
 f_{L,\varepsilon},\quad e_\varepsilon,\quad
                 \mathfrak b_Y f_{L,\varepsilon},\quad
                 \mathfrak b_Y e_\varepsilon
                            \in\mathcal U_0.
 \tag{MA-PI.92}
\]
In particular every theta image subsequently used is
in its actual all-exponential source image.

The next estimates are necessary: cutting off near zero
does not tend to zero in every Schwartz seminorm.
Because $f_L$ is even and $f_L(0)=0$, both its value and
first derivative vanish at zero. Therefore
\[
 r_\varepsilon(t)=\varepsilon^2g_\varepsilon(t/\varepsilon),
 \quad g_\varepsilon(y)=\omega(y)
                    f_L(\varepsilon y)/\varepsilon^2,
 \quad |c_1(\varepsilon)|+|c_2(\varepsilon)|=O(\varepsilon^3).
 \tag{MA-PI.93}
\]
The $g_\varepsilon$ form a bounded family of even compact
smooth functions with $g_\varepsilon(0)=0$. For derivatives
of order zero and one this follows from Taylor's formula
using the two vanishings; higher derivatives of the
non-cutoff factor have coefficients $\varepsilon^{j-2}$,
$j\ge2$. The last bound in (MA-PI.93) follows from
$\|r_\varepsilon\|_{L^1}=O(\varepsilon^3)$,
boundedness of $k_Y$, and the fixed inverse $T_Y^{-1}$.

\emph{5. A uniform actual theta estimate for the flattened error.}
Fix a smooth real step $\phi$ equal to one far left and
zero far right, and put $P(x)=\phi(x)e^{-x/2}$.
For every even Schwartz $v$ with $v(0)=0$, actual
Poisson summation gives
\[
 R_v(x):=\Theta v(x)-I(v)P(x)\in\mathcal E,
 \tag{MA-PI.94}
\]
continuously as a function of $v$ on that closed Schwartz
subspace. At the left end the exact Poisson formula is
$\Theta v(x)=I(v)e^{-x/2}+\Theta\widehat v(-x)$,
whose last term decreases faster than every exponential;
at the right end use the direct theta sum. The same
estimates apply to every derivative. This proves the
claimed continuity, including for bounded Schwartz
families, without applying theta continuity on $\mathcal U_0$
to a vector outside that subspace.

Set $q_\varepsilon=I(g_\varepsilon)$, a bounded scalar.
By (MA-PI.90)--(MA-PI.93),
$\sum_j c_jI(\psi_j)=I(r_\varepsilon)=\varepsilon^3q_\varepsilon$.
Scaling the actual theta series and subtracting its
polar part gives the exact identity
\[
 \begin{split}
 \Theta e_\varepsilon(x)
   ={}&\varepsilon^{5/2}R_{g_\varepsilon}(x-\log\varepsilon)
                      -\sum_{j=1}^2c_j(\varepsilon)R_{\psi_j}(x)\\
     &+\varepsilon^3q_\varepsilon e^{-x/2}
            [\phi(x-\log\varepsilon)-\phi(x)].
 \end{split}
 \tag{MA-PI.95}
\]
All terms on the right are in either specified $H$.
Put $c=a$ for the symmetric metric and $c=|b|$ for
the one-sided metric; thus $0\le c<1/2$ and the
one-sided norm is bounded above by the symmetric
$c$ norm. Translation of a bounded $\mathcal E$ family
by $-\log\varepsilon$ has norm at most
$C\varepsilon^{-c}(1+|\log\varepsilon|)^{s/2}$.
For the last term the bracket is supported in an
interval contained in $[\log\varepsilon-C,C]$.
Its absolute value is bounded independently of
$\varepsilon$. Hence
\[
 \begin{split}
 &\bigl\|\varepsilon^3e^{-x/2}
        [\phi(x-\log\varepsilon)-\phi(x)]\bigr\|_H^2\\
 &\qquad\le C\varepsilon^6
      \int_{\log\varepsilon-C}^{C}
           e^{-x+2c|x|}(1+x^2)^{s/2}dx\\
 &\qquad\le C'\varepsilon^{5-2c}
                          (1+|\log\varepsilon|)^s.
 \end{split}
 \tag{MA-PI.96}
\]
The fixed bounded positive part of the integral is
absorbed in the last bound. Combining (MA-PI.93)--(MA-PI.96)
proves
\[
 \|\Theta e_\varepsilon\|_H
   \le C\bigl[\varepsilon^{5/2-c}
                 (1+|\log\varepsilon|)^{s/2}+\varepsilon^3\bigr]
                                      \longrightarrow0.
 \tag{MA-PI.97}
\]
This estimate has kept the canceling polar pieces
together; none is separately presumed to be a vector
of the narrow Hilbert space.

There is also convergence after the fixed physical
cutoff. Direct differentiation of the Fourier integral gives
\[
 \sup_\xi|\partial_\xi^j\widehat r_\varepsilon(\xi)|
                         \le C_j\varepsilon^{j+3}.
 \tag{MA-PI.98}
\]
After multiplication by the fixed compact smooth
multiplier $\rho_Y(|\xi|)$, these estimates and
$c_j=O(\varepsilon^3)$ show that
$\mathfrak b_Y e_\varepsilon\to0$ in Schwartz space,
with every seminorm $O(\varepsilon^3)$. It belongs
to $\mathcal U_0$ by (MA-PI.92). The genuine continuous
theta map on that subspace therefore gives
\[
                 \|\Theta\mathfrak b_Ye_\varepsilon\|_H
                                      \longrightarrow0.
 \tag{MA-PI.99}
\]
The derivative constants may depend on the already
fixed profile; the flat approximation is chosen
after that profile and finite level have been fixed.

Finally use the bounded orthogonal projection $Q_Y$
and test against the actual vector $\Theta h_L\in J_Y$.
Equations (MA-PI.88), (MA-PI.97) and (MA-PI.99) give
\[
 \begin{split}
 &\langle Q_Y\Theta f_{L,\varepsilon}
            -\Theta\mathfrak b_Y f_{L,\varepsilon},\Theta h_L\rangle_H\\
 &\qquad\longrightarrow
                       \langle\Theta f_L,\Theta h_L\rangle_H\ne0.
 \end{split}
 \tag{MA-PI.100}
\]
Some finite positive $\varepsilon$ therefore supplies
$f_Y=f_{L,\varepsilon}$ satisfying (MA-PI.74).
Every input and output asserted there is an actual
Schwartz source vector with the required moments.
\end{proof}

\paragraph{The exact printed principal realization.}
The source isometry in (MA-PI.32)--(MA-PI.36) also has
the following direct trivial-character specialization.
Put $m=\mu_1(C_{\mathbb Q}^1)$ and $b=A-1/2$. For
$g=\nu(e^x)k$ and the source $f\otimes\mathbf1_{\widehat{\mathbb Z}}$,
\[
 \begin{aligned}
 (T_A^0f)(g)
   &:=|g|^A\Phi_{\mathbb Q}
                   (f\otimes\mathbf1_{\widehat{\mathbb Z}})(g)
       =e^{bx}\Theta f(x),\\
 (U_b^0F)(\nu(e^x)k)&=m^{-1/2}e^{bx}F(x).
 \end{aligned}
 \tag{MA-PI.101}
\]
The finite integer balls are unchanged by finite units,
so the theta function is constant on the compact fiber.
The two signs give the factor two in (MA-PI.70).
The exact printed norm (MA-PI.33) makes $U_b^0$ unitary
from $H^{\rm one}_{b,s}$ onto the trivial compact-character
component of $W^{s,2}(C_{\mathbb Q})$, and
$T_A^0=\sqrt m\,U_b^0\Theta$.
For $S=\{\infty\}$ the printed invariant PW source is
exactly the compact Fourier support condition in
(MA-PI.70). An arbitrary source can be replaced by its
even part without changing its theta sum or either
polar condition, and this preserves the Fourier band.
Therefore its actual printed orthogonal projection obeys
\[
 Q^{\rm src}_{\infty,Y,A,s}
                       =U_b^0Q_Y^{\rm one}(U_b^0)^{-1}.
 \tag{MA-PI.102}
\]
This is a unitary identity derived from the source
definitions, not a change to an equivalent metric.
For every $0<A<1$, (MA-PI.74) in its one-sided metric
consequently gives genuine flat witnesses with
\[
 Q^{\rm src}_{\infty,Y,A,s}T_A^0f_Y
                       \ne T_A^0\mathfrak b_Yf_Y.
 \tag{MA-PI.103}
\]
In particular this proves the same principal rational
canonical obstruction at the exact twist $A=1/2$
used on p.~515, with both actual source and cutoff
output polar free. The compact Haar mass has been
retained throughout the identification.

\paragraph{Source attribution and the precise conclusion.}
The theta source and rational PW support are the actual
MA-GP.1--2 objects. Their source locators are Masullo,
Definition 7.3.2.1, (7.3.2.1--7), p.~351, and
Definition 7.3.2.3, (7.3.2.12--17), pp.~354--355;
the Hilbert orthogonal projections are Definitions
11.2.1.1/11.2.1.3, pp.~487--489, with the explicit
metric here retained. The physical cutoff is the
first Step~II bullet on p.~516 and (12.4.1.8--10),
p.~517. Its action is computed on the actual source.
The even gamma factor agrees with Tate (1950),
Theorem~2.4.1 and the real table, thesis p.~(2.14),
PDF p.~21; its needed normalization and the precise
factorial consequence are derived above. These
are source interfaces, not a fresh audit of Tate's
whole global proof.

This proves failure of the canonical theta-map
intertwining $Q_Y\Theta=\Theta\mathfrak b_Y$ in the
main rational trivial sector, on the corrected
symmetric metric and separately on each displayed
one-sided metric. It does not assert that the draft
prints that exact equation, or exclude every different
nonlocal or Gram-corrected comparison map. The actual
Gram factors of MA-GP remain valid. The scalar trace
identity supplied by $A$, the independent FT002
deficiency gate, and the distinction between a
restricted ideal return and the full quotient trace
are unchanged.

\paragraph{Grothendieck placement.}
The arithmetic sector is the trivial rank-one Artin
representation, the constant finite-monodromy local
system on $\operatorname{Spec}\mathbb Z$, with Euler
factor $(1-p^{-z})^{-1}$. Its placement is SGA~1,
Expos\'e~V, \S7, and Deligne, \emph{Valeurs de fonctions L},
Definition~6.1 and \S\S6.2--6.3, in the inherited
conventions of MA-CH and MA-CQ. The two source moments
are analytic polar conditions. The metric obstruction
does not change the Artin representation, a $K_0$
class, or a motivic weight. In cote~1, scans~079--080
and~111--112 through BR-TR-01/02, the trace/determinant
comparison still requires defined continuous/nuclear
maps. This proof refutes one exact proposed map while
leaving the previously constructed actual comparison
operators at their proved scopes.
